% concatenated from Submission-v2.tex
%%%
%%%             Submission-v2
%%%


\documentclass[11pt]{article}
%\documentclass[a4paper]{report}  
\usepackage{mathrsfs} 
\usepackage[margin=1.25in]{geometry}
\usepackage{ifthen}
\usepackage{graphicx}
\usepackage{epsfig}
\usepackage{dsfont, amssymb}
\usepackage{amsfonts,dsfont,amssymb,amsthm,stmaryrd,bbm}
\usepackage[cmex10]{amsmath} % Use the [cmex10] option 
   % to ensure compliance with IEEE Xplore (see bare_conf.tex)
\usepackage{soul}

%\usepackage[T1]{fontenc}
%\usepackage[applemac]{inputenc}		
\usepackage{hyperref,mathtools}
\hypersetup{colorlinks=true,pdftitle=""}

\setlength{\topmargin}{-0.45 in}     %% dist from automatic 1in baseline to top of header
\setlength{\oddsidemargin}{0.3in}  %% automatic 1in baseline
\setlength{\evensidemargin}{0.3in} 
\setlength{\textheight}{9in}
\setlength{\textwidth}{6.1in} 
\setlength{\footskip}{0.55in}  %% dist(footer bo




%\theoremstyle{definition}
%\newtheorem{definition}{Definition}%[section]
%\theoremstyle{plain}
%\newtheorem{thm}{Theorem}%[section]
%\theoremstyle{plain}
%\newtheorem{prop}{Proposition}%[section]
%\theoremstyle{plain}
%\newtheorem{lemma}{Lemma}%[section]
%\theoremstyle{plain}
%\newtheorem{corol}{Corollary}%[thm]
%\theoremstyle{plain}
%\newtheorem{conjecture}{Conjecture}%[section]
%\theoremstyle{remark}
%\newtheorem{remark}{Remark}%[section]
%\theoremstyle{plain}
%\newtheorem{conj}{Conjecture}%[section]


%\newtheorem{conj}{Conjecture}[section]
\newtheorem{thm}{Theorem}[section]
%\newtheorem{ce}[conj]{Counter-example}
\newtheorem{remark}[thm]{Remark}
\newtheorem{lem}[thm]{Lemma}
\newtheorem{prop}[thm]{Proposition}
\newtheorem{coro}[thm]{Corollary}
%\newtheorem{ques}[conj]{Question}
%\newtheorem{cond}[conj]{Condition}
\newtheorem{defn}[thm]{Definition}
\newtheorem{cor}[thm]{Corollary}
%\newtheorem{ex}[conj]{Example}

\newcommand\independent{\protect\mathpalette{\protect\independent}{\perp}} 
\def\independent#1#2{\mathrel{\rlap{$#1#2$}\mkern2mu{#1#2}}} 


\newcommand{\supp}{\mathrm{Supp}} 
\newcommand{\cF}{\mathcal{F}} 
%\newcommand{\cE}{\mathcal{E}} 
\newcommand{\F}{\mathbb{F}}  
\newcommand{\mR}{\mathbb{R}} 
\newcommand{\R}{\mathbb{R}} 
\newcommand{\mN}{\mathbb{N}}
\newcommand{\mZ}{\mathbb{Z}}
\newcommand{\mC}{\mathbb{C}}
\newcommand{\mP}{\mathbb{P}}
\newcommand{\mF}{\mathbb{F}}
\newcommand{\mE}{\mathbb{E}}
\newcommand{\mS}{\mathbb{S}}
\newcommand{\ima}{\mathrm{Im}\:}
\newcommand{\re}{\mathrm{Re}\:}
\newcommand{\tr}{\mathrm{tr}}
\newcommand{\Corr}{\mathrm{corr}}
\newcommand{\Cov}{\mathrm{cov}}
\newcommand{\pp}{\mathbb{P}}
\DeclareMathOperator{\Var}{Var}
\newcommand{\e}{\varepsilon}
%\newcommand{\lam}{\lambda}
\newcommand{\ent}{h}
\newcommand{\info}{I}
\newcommand{\rank}{\mathrm{rank}}
\DeclareMathOperator{\diag}{diag}
\newcommand{\A}{\mathcal{A}}
\newcommand{\X}{\mathcal{X}}
\newcommand{\Y}{\mathcal{Y}}
%\newcommand{\Z}{\mathcal{Z}}
\newcommand{\jj}{[j]}
\newcommand{\ii}{[i]}
\newcommand{\ic}{[i^c]}
\newcommand{\J}{[J]}
\newcommand{\C}{[J^c]}
\newcommand{\K}{[J\setminus i]}
\newcommand{\m}{[m]}
\newcommand{\1}{[1]}
\newcommand{\M}{E_m}
\newcommand{\bX}{\bar{X}}
\newcommand{\bY}{\bar{Y}}
\newcommand{\bU}{\bar{U}}
\newcommand{\mat}{\mathrm{MAT}}
\newcommand{\JJ}{\hat{I}}

%%%%%%%%%%%%


\newcommand{\de}{\mathrm{d}}
%\newcommand{\e}{\mathrm{e}}
\newcommand{\conv}{\mathrm{conv}}
\newcommand{\relint}{\mathrm{relint}}
\newcommand{\vol}{\mathrm{Vol}}
\newcommand{\diam}{\mathrm{diam}}
\newcommand{\card}{\mathrm{card}}
\newcommand{\interior}{\mathrm{int}}
\newcommand{\proj}{\mathrm{P}}
\newcommand{\vect}{\mathrm{vect}}
%\newcommand{\supp}{\mathrm{supp}}
\newcommand{\dist}{\mathrm{dist}}
\newcommand{\Hess}{\mathrm{Hess}}

%\newcommand{\R}{\mathbb{R}}
\newcommand{\N}{\mathbb{N}}
\newcommand{\Z}{\mathbb{Z}}
%\newcommand{\F}{\mathcal{F}}





%\newcommand{\setS}{\mbox{${\bf s}$}}
\newcommand{\var}{\VAR}
\newcommand{\setS}{s}
\newcommand{\setT}{t}
\newcommand{\setU}{u}
\newcommand{\nullset}{\phi}
\newcommand{\Xs}{X_{\setS}}
\newcommand{\gs}{\gamma_{\setS}}
\newcommand{\collS}{\mathcal{C}}
\newcommand{\Etb}{\bar{E}_{\setT}}
\newcommand{\asb}{\bar{a}_{\setS}}
\newcommand{\ws}{w_{\setS}}
\newcommand{\mth}{\frac{1}{m}}
\newcommand{\rth}{\frac{1}{r_{+}}}
\newcommand{\sumS}{\sum_{\setS\in\collS}} 
\newcommand{\prodS}{\prod_{\setS\in\collS}} 
\newcommand{\mus}{\mu_{\setS}}
\newcommand{\as}{\alpha_{\setS}}
\newcommand{\lams}{\lam_{\setS}}
\newcommand{\bs}{\beta_{\setS}}
\newcommand{\Zs}{Z_{\lams}}
\newcommand{\dth}{\frac{1}{d}}
\newcommand{\calN}{\mathcal{N}}
\newcommand{\Prob}{\mathbb{P}}
\newcommand{\eqq}{\mbox{$ \;\stackrel{(?)}{=}\; $}}


\def\E{{\bf E}}
%\def\R{{\bf R}}
\def\C{{\bf C}}
\def\P{{\bf P}}
\def\H{{\cal H}}
\def\Im{{\rm Im}}
\def\Tr{{\rm Tr}}
%\newcommand{\Hess}{\text{Hess}}

\def\k{{\kappa}}
\def\M{{\cal M}}
\def\Var{{\rm Var}}
\def\Ent{{\rm Ent}}
\def\O{{\rm Osc}_\mu}

\def\ep{\varepsilon}
\def\phi{\varphi}
\def\F{{\cal F}}
\def\L{{\cal L}}

\def\bee{\begin{eqnarray*}}
\def\ene{\end{eqnarray*}}
\usepackage{mathtools}
\DeclarePairedDelimiter\ceil{\lceil}{\rceil}
\DeclarePairedDelimiter\floor{\lfloor}{\rfloor}

%%%%%%%%





\title{A note on the Littlewood-Offord problem for discrete log-concave distributions}
\author{Arnaud Marsiglietti\footnote{Department of Mathematics, University of Florida, Gainesville, FL 32611, USA, E-mail: a.marsiglietti@ufl.edu} \, and James Melbourne\footnote{Probabilidad y Estad\'isticas, Centro de Investigaci\'ones en Matem\'aticas, Guanajuato, GTO 36023, MX, E-mail: james.melbourne@cimat.mx}}


\date{}



\begin{document}

\maketitle

\begin{abstract}

We present an extension of the famous Littlewood-Offord problem when Bernoulli distributions are replaced with discrete log-concave distributions. A variant of the Littlewood-Offord problem for arithmetic progressions, as well as an entropic version, is also discussed. Along the way, we recover and extend a result of Madiman and Woo (2015) on the entropy power inequality for discrete uniform distributions.

\end{abstract}

\vskip5mm
{\bf Keywords:} Littlewood-Offord problem, log-concave, arithmetic progression, R\'enyi entropy, majorization


%\vskip1cm
\section{Introduction}

Given $a = (a_1, \dots, a_n) \in (\R \setminus \{0\})^n$ and independent random variables $X_1, \dots, X_n$, $n \geq 1$, with a Rademacher distribution, that is, for all $1 \leq k \leq n$, $\Prob(X_k= \pm 1) = 1/2$,
the question of estimating the quantity
$$ \sup_{x \in \R} \, \Prob(a_1 X_1 + \cdots + a_n X_n = x) $$
is often referred to as the Littlewood-Offord problem. A classical result of Littlewood and Offord \cite{LO} and Erd\"os \cite{E} states that
\begin{eqnarray}\label{erdos}
\sup_{a \in (\R \setminus \{0\})^n} \sup_{x \in \R} \, \Prob(a_1 X_1 + \cdots + a_n X_n = x) \leq \frac{1}{2^n}\binom{n}{\lfloor \frac{n}{2} \rfloor} = O \bigg( \frac{1}{\sqrt{n}} \bigg).
\end{eqnarray}

Kleitman extended this result when $a_1, \dots, a_n$ are vectors in a Hilbert space \cite{K1}, \cite{K2}. Many variants of the Erd\"os-Littlewood-Offord problem have been established, such as improved bounds under certain constraints on the $a_i$'s \cite{H}, an inverse Littlewood-Offord theorem \cite{TV}, \cite{NV}, and a resilience version \cite{BFK}, but most are mainly dealing with Bernoulli distribution with parameter $1/2$. Recently, Fox, Kwan and Sauermann \cite[Question 6.2]{FKS} asked whether the Littlewood-Offord problem can be solved for Bernoulli distribution of arbitrary parameter $p \in (0,1)$. This question has been investigated by Singhal \cite{S}, who gave a qualitative solution to the problem, showing that a maximizer is obtained for $a_i \in \{-1, 1\}$ and gave Fourier theoretic bounds. Sharp quantitative bounds were found by Madiman, Melbourne, and Roberto in \cite{MMR} where an entropic generalization of the problem was considered. The case of general random variables was considered recently by Juškevičius and Kurauskas \cite{JK}.

Recall that an integer-valued random variable $X$ is said to be discrete log-concave if its probability mass function $p$ satisfies
$$ p(j)^{2} \geq p(j-1)p(j+1) $$
for all $j \in \mathbb{Z}$ and the support of $X$ is an integer interval. Discrete log-concave distributions form an important class. Examples include discrete uniform, Bernoulli, binomial
and convolutions of Bernoulli distributions with arbitrary parameters, Poisson, geometric, negative binomial, etc. (cf. \cite{JG} and references therein). We refer to \cite{Sta}, \cite{Bre}, \cite{SW}, \cite{Bra}  for further background on log-concavity.

The goal of this article is to extend the solution of the Littlewood-Offord problem to the whole class of discrete log-concave distributions. Our bounds are quantitative and non-asymptotic. In particular, we prove the following.


\begin{thm}\label{main1}

Let $X_1, \dots, X_n$, $n \geq 1$, be independent discrete log-concave random variables finitely supported. Then,
\begin{equation}\label{1}
\sup_{a \in (\R \setminus \{0\})^n} \sup_{x \in \R} \, \Prob(a_1 X_1 + \cdots + a_n X_n = x) \leq \frac{1}{\sqrt{1 + c \sum_{k=1}^n \Var(X_k)}},
\end{equation}
with $c=1$. Moreover, one may take $c = 2$ when the random variables are, in addition, symmetric about a point.

\end{thm}

We note that this recovers the $O(1/\sqrt{n})$ bound for independent Bernoulli distribution with parameter $1/2$.

It has been shown in \cite{MMR} that the bound \eqref{1} holds with $c=2$ when the $X_k$'s have a Bernoulli distribution with arbitrary parameters.

Theorem \ref{main1} is sharp up to an absolute constant, as the left-hand side of inequality \eqref{1} can be lower bounded by
$$ \frac{1}{\sqrt{1 + 12\sum_{k=1}^n \Var(X_k)}}, $$
see Remark \ref{sharp}.


The article also presents an entropic version of the Littlewood-Offord problem. See Section~\ref{2} for a precise definition of the R\'enyi entropy power $N_{\alpha}$.

\begin{thm}\label{main2}
Let $\alpha \in [0, +\infty]$ and $n \geq 1$. Let $X_1, \dots, X_n$ be independent discrete log-concave random variables finitely supported. Then,
\begin{equation}\label{entrop}
\inf_{a \in (\R \setminus \{0\})^n} N_\alpha(a_1 X_1 + \cdots + a_n X_n) \geq 1 + c \sum_{k=1}^n \Var(X_k),
\end{equation}
with $c=1$. Moreover, one may take $c = 4$ when $1 < \alpha \leq 2$, and for other values of $\alpha$ one may take $c=2$ when the random variables are, in addition, symmetric about a point.
\end{thm}

It has been shown in \cite{MMR} that the bound \eqref{entrop} holds with $c = \frac{2 \alpha}{\alpha - 1}$, $\alpha \geq 2$, when the $X_k$'s have a Bernoulli distribution. It turns out that Theorem \ref{main1} is a particular case of Theorem \ref{main2}.

We also present a version of the Littlewood-Offord problem for arithmetic progressions.

\begin{thm}\label{main3}
Let $X_1, \dots, X_n$, $n \geq 1$, be independent discrete log-concave random variables finitely supported. Then,
\begin{equation*}
\sup_{a \in (\R \setminus \{0\})^n} \sup_{x \in \R} \, \Prob(a_1 X_1 + \cdots + a_n X_n \in A_{l,m}(x)) \leq \frac{l}{\sqrt{1 + c \sum_{k=1}^n \Var(X_k) + c \, \frac{l^2-1}{12}}},
\end{equation*}
with $c=1$. Moreover, one may take $c = 2$ when the random variables are, in addition, symmetric about a point.
\end{thm}
Here, $A_{l,m}(x)$ is an arithmetic progression of length $l \geq 1$, that is $A_{l,m}(x) = \{x + mj \}_{j=1}^l$ for $m \in \mathbb{Z}$ and $x \in \mathbb{R}$. In fact, $m$ can be taken as a real number (see Section \ref{4}). For example, if the $X_k$'s are i.i.d. Bernoulli with parameter $1/2$, then we deduce
\[
    \sup_{a \in (\R \setminus \{0\})^n} \sup_{x \in \R} \, \Prob(a_1 X_1 + \cdots + a_n X_n \in A_{l,m}(x)) \leq \frac l {\sqrt{ 1 + \frac n 2 + \frac{l^2-1}{6}}}.
\]
We refer to Section \ref{4} for an estimate that improves Theorem \ref{main3} when the $X_k$'s have a Bernoulli distribution with arbitrary parameter $p \in (0,1)$, $p \neq \frac{1}{2}$.  

Let us note that the case $l = 1$ corresponds to the classical Littlewood-Offord problem, hence Theorem \ref{main1} is a particular case of Theorem \ref{main3}.

Finally, our method allows us to establish the following entropy power inequality for discrete uniform distributions.

\begin{thm}\label{EPI-unif}
    Let $\alpha \in [0, 2]$ and $n \geq 1$. Let $U_1, \dots, U_n$ be uniformly distributed independent random variables on any set of integers. Then,
    $$ N_{\alpha} \left( \sum_{k=1}^n U_k \right) \geq \sum_{k=1}^n N_{\alpha}(U_k) - (n-1). $$
\end{thm}

The particular case $\alpha = 1$ and $n=2$ of Theorem \ref{EPI-unif} has been established by Madiman and Woo in \cite{MW}.

The article is organized as follows. Section \ref{prel} provides the necessary background on the notion of majorization and rearrangement inequalities, which are the main ingredients in the proofs of Theorems \ref{main1} and \ref{main2}. The proofs of Theorems \ref{main1} and \ref{main2} are postponed to Section~\ref{2}. Section \ref{sec-4} demonstrates the advantage of our results compared with general bounds on the concentration function existing in the literature. Section \ref{4} presents a variant of the Littlewood-Offord problem for arithmetic progressions, in particular, Theorem \ref{main3} is proved. The last section focuses on the Bernoulli and uniform distributions, where Theorem \ref{EPI-unif} is proved.




\section{Preliminaries}\label{prel}

Throughout the article, we use the notation $x \cdot y = \sum_{i=1}^n x_i y_i$, for the dot product of $x=(x_1, \dots, x_n)$ and $y=(y_1, \dots, y_n)$. We also denote
\begin{equation}\label{M}
M(X) = \sup_{x \in \R} \, \Prob(X = x).
\end{equation}


We will need the following result.


\begin{thm}[\cite{A}, \cite{BMM}]\label{LC}

If the random variable $X$ has a discrete log-concave distribution, then
\begin{equation}\label{M-bound}
\frac{1}{\sqrt{1 + 12\Var(X)}} \leq M(X) \leq \frac{1}{\sqrt{1 + \Var(X)}}.
\end{equation}
Moreover, if the distribution of $X$ is symmetric about a point, then the upper bound may be sharpened to
$$ M(X) \leq \frac{1}{\sqrt{1 + 2\Var(X)}}. $$

\end{thm}

Let us note that the lower bound in \eqref{M-bound} holds for arbitrary random variables. The upper bound in \eqref{M-bound} was proven by Aravinda \cite{A}, who refined the bound
$$ M(X) \leq \frac{2}{\sqrt{1 + 4\Var(X)}} $$
obtained by Bobkov and the authors in \cite{BMM} (see also \cite{GMPP}, \cite{JMNS}).



Recall that a probability distribution $p = (p_1, \dots, p_n)$ written in descending order with positive entries, $p_i \geq p_{i+1}$, is majorized by another $q = (q_1, \dots, q_m)$ also written in descending order with positive entries when 
\begin{align}
    \sum_{i=1}^k q_i \geq \sum_{i=1}^k p_i
\end{align}
holds for all $k$. We write $p \prec q$ when $p$ is majorized by $q$. We also write $X \prec Y$ when the probability mass function of $X$ is majorized by the probability mass function of $Y$. After extending $q$ to $\{1, \dots, n\}$ by setting $q_j = 0$ for $m+1 \leq j \leq n$, this is equivalent to the existence of a doubly stochastic matrix $A$ such that $Aq = p$ (see \cite{MOA}). 

For example, if $q$ is a point mass, $(q_1, \dots, q_n) = (1,0,\dots,0)$, then for any other $p = (p_1, \dots, p_n)$ we can write
\begin{align}
    A = \left( \begin{array}{cccc}
    p_1 &p_2 &\dots &p_n \\
    p_2 &p_3 & \dots &p_{1} \\
    \vdots &\vdots & \vdots & \vdots \\
    p_n & p_1 & \dots & p_{n-1}
    \end{array} 
    \right),
\end{align}
so that $Aq = p$. More generally, if $q$ is not a probability sequence and $q = (M,0,\dots, 0)$ while $\sum_i p_i = M$ then setting $\lambda_i = p_i/M$, we can write
\begin{align} \label{eq: double stochastic matrix for point mass}
    A = \left( \begin{array}{cccc}
    \lambda_1 &\lambda_2 &\dots &\lambda_n \\
    \lambda_2 &\lambda_3 & \dots &\lambda_{1} \\
    \vdots &\vdots & \vdots & \vdots \\
    \lambda_n & \lambda_1 & \dots & \lambda_{n-1}
    \end{array} 
    \right),
\end{align}
so that $Aq = p$.


\begin{lem} \label{lem: majorization lemma}
If $Y$ is a random variable taking finitely many values, and $f$ is a deterministic function, then $Y \prec f(Y)$.
\end{lem}


\begin{proof}
Without loss of generality, we may assume that $Y$ is distributed on $\{1,2,\dots, n\}$, and denote by $\{a_1, \dots, a_m\}$ the support of $f(Y)$. If the distribution of $Y$ is denoted by $p$, note that the distribution of $f(Y)$, written as $q$, will satisfy $q_i = \sum_{j \in f^{-1}(\{a_i\})} p_j$.  Writing $q = (q_1,0, \dots,0,q_2, 0, \dots,0, \dots, q_m, 0,\dots, 0)$ where the number of zeros between $q_i$ and $q_{i+1}$ is determined by the cardinality $n_i$ of $f^{-1}(\{a_i\})$, and writing $p$ in order such that $f^{-1}(\{a_i\}) = \{k_i, k_i+1, \dots, k_i+n_i-1\}$ so that $q_i = p_{k_i} + \cdots + p_{k_i + n_i - 1}$, based on the discussion above we can write a doubly stochastic block matrix,
\begin{align}
    \Lambda = \left( \begin{array}{ccccc}
        (A_1)&(0) &(0) &\dots & (0) \\
        (0) & (A_2) &(0) &\dots & (0) \\
        \vdots & \vdots & \ddots && \vdots \\
        \vdots & \vdots & & \ddots & \vdots \\
        (0) & (0) & \dots & (0) & (A_m)
        \end{array}
        \right)
\end{align}
where each $A_i$ is of the form of $\eqref{eq: double stochastic matrix for point mass}$ for the $p_j$ such that $f(j) = a_i$, so that $\Lambda q  = p$, and the lemma holds.
\end{proof}

If $f \colon \Z \to [0, +\infty)$ is finitely supported, with support $\{x_0, \dots, x_m\}$, denote by $f^\#$ its squeezed rearrangement, that is, $f^\#$ is supported on $\{0, \dots,m\}$ and $f^\#(j) = f(x_j)$, for $j\in\{0, \dots,m\}$. If $X$ is an integer-valued random variable with probability mass function $f$, we denote by $X^\#$ the random variable with probability mass function $f^\#$. The following result was proven in \cite{MWW}.

\begin{thm}[\cite{MWW}]\label{rearrance-2}
If $X_1, \dots, X_n$ are integer-valued independent random variables such that $X_1^\#, \dots, X_n^\#$ are log-concave, then
$$ X_1 + \dots + X_n \prec X_1^\# + \dots + X_n^{\#}. $$    
\end{thm}




Finally, recall that the R\'enyi entropy of order $\alpha \in (0,1) \cup (1, +\infty)$ of a discrete random variable $X$ with values in a countable set $I$ and with probability mass function $p$ is defined as
$$ H_{\alpha}(X) = \frac{1}{1-\alpha} \log \left( \sum_{x \in I} p^{\alpha}(x) \right) = \log(\|p\|_{\alpha}^{\frac{\alpha}{1-\alpha}}). $$
The limit cases are interpreted as
$$ H_0(X) = \log(|\supp(p)|), \quad H_1(X) = -\sum_{x \in I} p(x) \log(p(x)), \quad H_{\infty}(X) = -\log(\sup_{x \in I} p(x)). $$
We note that the $M$-functional defined in \eqref{M} may be viewed as a member of the family of R\'enyi entropies via the formula
\begin{equation}\label{member}
M(X) = e^{-H_{\infty}(X)}.
\end{equation}
In particular, considering the R\'enyi entropy power $N_{\alpha}(X) = e^{2H_{\alpha}(X)}$, Theorem \ref{LC} yields the bound
\begin{equation}\label{entropy}
N_{\alpha}(X) \geq N_{\infty}(X) \geq 1 + \Var(X),
\end{equation}
for arbitrary discrete log-concave random variable $X$, where the first inequality holds by monotonicity of R\'enyi entropy. 




\section{The Littlewood-Offord problem for discrete log-concave distributions and an entropic version}\label{2}

Throughout this section, given random variables $X_1, \dots, X_n$, we denote $X=(X_1, \dots,X_n)$. The first step in establishing the Littlewood-Offord problem for log-concave distributions is to reduce the problem to signs.


\begin{thm}\label{thm: majorization by signs}

For $a_i \in \mathbb{R} \setminus \{0\}$, and $X_i$ independent, log-concave, $\mathbb{Z}$-valued random variables taking finitely many values, there exist signs $v_i \in \pm 1$ such that $a \cdot X \prec v \cdot X$.

\end{thm}


\begin{proof}
Let us first observe that $a \cdot X$ can be majorized by $\tilde{a} \cdot X$, for set of constants $\tilde{a}_i \in \mathbb{Z} \setminus \{0\}$. Indeed, it has been observed in \cite[Proof of Lemma 5.1]{MMR} that one may construct a linear map $T \colon \R \to \mathbb{Q}$ such that $T(a_i) \in \mathbb{Z} \setminus \{0\}$ for all $i$. Thus, by Lemma \ref{lem: majorization lemma}, $a \cdot X \prec T( a\cdot X) $. Further, since the $X_i$'s are integer-valued, one has
\begin{align}
T(a_1 X_1 + \cdots + a_n X_n)
    = T(a_1) X_1 + \cdots + T(a_n) X_n.
\end{align}
Writing $T(a) = (T(a_1), \dots, T(a_n)) \in (\mathbb{Z} \setminus \{0\} )^n$, we thus have $a \cdot X \prec T(a) \cdot X$.

Observe that our result follows from Theorem \ref{rearrance-2} since $T(a) \cdot X \prec (T(a_1) X_1)^{\#} + \cdots + (T(a_n) X_n )^{\#}$. Indeed, $(T(a_i) X_i)^{\#} = v_i X_i$ where $v_i = sign(T(a_i))$. Since $v_i X_i$ is log-concave one may apply Theorem \ref{rearrance-2} and we have $a \cdot X \prec T(a) \cdot X \prec v \cdot X$ where $v_i = \pm 1$ and the result follows.
\end{proof}


Since $\alpha$-R\'enyi entropy is Schur concave as a consequence of \cite[Proposition 3-C.1]{MOA} (see also \cite{MWW}),
%(Appendix F.3.a (p. 562)
the proof of Theorem \ref{main2} follows immediately.

\begin{proof}[Proof of Theorem \ref{main2}]
Theorem \ref{thm: majorization by signs} combined with Schur concavity of R\'enyi entropy yields that for $\alpha \in [0,\infty]$, $X = (X_1, \dots, X_n)$ with $X_i$ independent and log-concave  then for all $a = (a_1, \dots, a_n)$ with $a_i \in \mathbb{R} \setminus \{0\}$, there exists $v_i = \pm 1$ such that $v = (v_1, \dots, v_n)$ implies
\begin{align}
N_\alpha( a \cdot X) \geq N_\alpha( v \cdot X).
\end{align}
Moreover the choice of signs is independent of the $X_i$ determined only by the coefficients $a_i$.
Therefore, using \eqref{entropy},
$$ N_\alpha( a \cdot X) \geq N_\alpha( v \cdot X) \geq 1 +\Var(v \cdot X) = 1 + \sum_{i=1}^n \Var(X_i). $$
It has been shown in \cite{BMM} that when $1 < \alpha \leq 2$, $N_{\alpha}(X) \geq 1 + 4\Var(X)$ for arbitrary log-concave $X$, while $N_{\alpha}(X) \geq 1 + 2\Var(X)$ when $X$ is symmetric about a point and log-concave. This concludes the proof.    
\end{proof}


\begin{remark}\label{sharp}

Theorem \ref{main2} is sharp up to an absolute constant. This is a consequence of the bound $N_{\alpha}(X) \leq 1 + \frac{4(3 \alpha - 1)}{\alpha - 1} \Var(X)$ proved in \cite{BMM}, which holds for all $\alpha > 1$. Therefore,
$$ \inf_{a \in (\R \setminus \{0\})^n} N_\alpha(a \cdot X) \leq 1 + \frac{4(3 \alpha - 1)}{\alpha - 1} \sum_{i=1}^n \Var(X_i). $$
To obtain a better estimate when $\alpha$ tends to 1, one may use the following well-known upper bound for the discrete entropy, $N(X) \leq \frac{2 \pi e}{12} + 2\pi e \Var(X)$ (see, e.g., \cite{Massey}, \cite{BM}), which yields for all $\alpha \geq 1$,
$$ \inf_{a \in (\R \setminus \{0\})^n} N_\alpha(a \cdot X) \leq \frac{2 \pi e}{12} + 2 \pi e\sum_{i=1}^n \Var(X_i). $$
    
\end{remark}


Specializing Theorem \ref{main2} to $\alpha = +\infty$, and recalling \eqref{member}, we obtain Theorem \ref{main1}:
\begin{equation}\label{M-function}
\sup_{a \in (\R \setminus \{0\})^n} M(a \cdot X) \leq \frac{1}{\sqrt{1 + \sum_{i=1}^n \Var(X_i)}},
\end{equation}
holding for arbitrary independent log-concave random variables $X_1, \dots, X_n$. In particular, we deduce the following.










\begin{prop}\label{answer}

Let $X_1, \dots, X_n$ be i.i.d. Bernoulli distribution of parameter $p$. Then,
$$ \sup_{a \in (\R \setminus \{0\})^n} M(a \cdot X) \leq \frac{1}{\sqrt{1 + np(1-p)}}. $$

\end{prop}



As mentioned in the introduction, in this specific case of Bernoulli distribution, a refined argument has been used in \cite{MMR} to show that 
\begin{equation}\label{answer-2}
\sup_{a \in (\R \setminus \{0\})^n} M(a \cdot X) \leq \frac{1}{\sqrt{1 + 2np(1-p)}}.
\end{equation}


We note that one may provide a unification of both Erdos' result of the Littlewood-Offord problem \eqref{erdos} and \eqref{answer-2}.



\begin{prop}\label{extend}

Let $X_1, \dots, X_n$ be independent random variables such that for each $i \in \{1, \dots, n\}$, $X_i \in \{x_i, x_{i+1}\}$, where $x_i, x_{i+1} \in \Z$ with $x_i \leq x_{i+1}$ and 
$$ \Prob(X_i = x_i) = 1 - \Prob(X_i = x_{i+1}) = 1-\theta_i, \qquad \theta_i \in (0,1). $$
Then,
$$ \sup_{a \in (\R \setminus \{0\})^n} M(a \cdot X) \leq \frac{1}{\sqrt{1 + 2\sum_{i=1}^n \frac{\Var(X_i)}{(x_i-x_{i+1})^2}}}. $$

\end{prop}

\begin{proof}
Using the same argument as in the proof of Theorem \ref{thm: majorization by signs}, one may deduce that for all $a \in (\R \setminus \{0\})^n$,
$$ M(a \cdot X) \leq M(\mathbbm{1} \cdot B), $$
where $B_1, \dots, B_n$ are independent Bernoulli distributions of parameter $\theta_i$ or $1-\theta_i$ depending on the sign of $a_i$.
%assume that the $a_i$'s are integer valued. Denote $Y_i = a_i X_i$, and note that $Y_i^{\#}$ is a Bernoulli distribution of parameter $\theta_i$ if $a_i \geq 0$ and of parameter $1-\theta_i$ if $a_i \leq 0$. 
%In any case, $\Var(B_i) = $.
The result follows by using the bound
$$ M(\mathbbm{1} \cdot B) \leq \frac{1}{\sqrt{1 + 2\sum_{i=1}^n \Var(B_i)}}, $$
which is a consequence of \eqref{answer-2}, and by noting that
$$ \Var(B_i) = \theta_i(1-\theta_i) = \frac{\Var(X_i)}{(x_{i+1} - x_i)^2}. $$
\end{proof}


The Littlewood-Offord solution \eqref{erdos}, as well as Proposition \ref{answer} immediately follow from Proposition \ref{extend}. Note that one may even allow Rademacher $\pm 1$ distributions with arbitrary parameter $p \in (0,1)$, which yields the same inequality as for Bernoulli distributions. We state this result in the next corollary.

\begin{coro}

If the $X_i$'s are independent Rademacher distributions with arbitrary parameter $p \in (0,1)$, then
$$ \sup_{a \in (\R \setminus \{0\})^n} M(a \cdot X) \leq \frac{1}{\sqrt{1 + 2np(1-p)}}. $$

\end{coro}






\section{Comparison with general bounds on the concentration function}\label{sec-4}


The goal of this section is to demonstrate that our upper bound on
$$ M(X) = \sup_{x \in \mathbb{R}} \mathbb{P}(X = x) $$
given in equation \eqref{M-function}, specialized to $a=(1, \dots, 1)$, provides better information compared with existing results in the literature on concentration function. Let us recall that the concentration function of a random variable $X$ is defined as
$$ Q(X ; \lambda) = \sup_x \Prob(x \leq X \leq x + \lambda), \qquad \lambda \geq 0. $$
Note that for $X$ integer valued and $\lambda < 1$,
\[
    Q(X; \lambda) = M(X).
\]
A general bound was established by Miroshnikov and Rogozin \cite{M-R}.


\begin{thm}[Miroshnikov-Rogozin \cite{M-R}]
    There exists a universal constant $C > 0$ such that $2 \lambda \geq \lambda_k$ and $X_k$ independent random variables with $S = X_1 + \cdots + X_n$,
    \[
        Q(S; \lambda) \leq C \lambda \left( \sum_{k=1}^n \mathbb{E}\left( |X_k^s| \wedge \frac{\lambda_k}{2} \right)^2 Q^{-2}(X_k, \lambda_k) \right)^{-\frac 1 2},
    \]
    where given $X$, $X^s = X - X'$ and $X'$ is an independent copy of $X$.
\end{thm}

When $X_k \sim$ Bernoulli($p_k$), note that $|X_k^s| \sim$ Bernoulli$(2p_k(1-p_k))$ so that for $\lambda < 1$,
\[
    M(S) = Q(S;\lambda) \leq C \lambda \left( \sum_{k=1}^n  \frac{\lambda_k^2 \,2 p_k(1-p_k) }{4}  M^{-2}(X_k) \right)^{-\frac 1 2}.
\]
Minimizing the right hand side with $\lambda_k = 2 \lambda$ we have
\begin{equation}\label{MR}
    M(S) \leq C \left( \sum_{k=1}^n \frac{2 p_k (1-p_k)}{(p_k \vee (1-p_k))^2}\right)^{- \frac 1 2} = C \left( \sum_{k=1}^n \frac{ 2 \Var(X_k)}{(p_k \vee (1-p_k))^2} \right)^{- \frac 1 2}.
\end{equation}
\iffalse
Since $\frac 1 2 \leq p \vee (1-p) \leq 1$ for $p \in [0,1]$, we have that for $\lambda < 1$ 
\begin{equation}\label{MR}
    Q(S;\lambda) = M(S) \leq \frac{C}{\sqrt{2}} \left( \sum_{k=1}^n \Var(X_k) \right)^{-\frac 1 2} \leq  2 C \lambda \left( \sum_{k=1}^n \mathbb{E} \left( |X_k^s | \wedge \frac{\lambda_k}{2} \right)^2 Q^{-2}(X_k, \lambda_k) \right)^{- \frac 1 2}  
\end{equation}
for any choice of $\lambda_k \leq 2 \lambda$.
\fi
Since $\frac 1 2 \leq p \vee (1-p) \leq 1$ for $p \in [0,1]$, the right-hand side of \eqref{MR} is at least $\frac{C}{\sqrt{2}} \left( \sum_{k=1}^n \Var(X_k) \right)^{-\frac 1 2}$. Therefore for Bernoulli distributions, even after optimizing, the Miroshnikov-Rogozin yields the bound
$$ M(S) \leq \frac{\widetilde{C}}{\sqrt{\Var(S)}} $$
for some absolute constant $\widetilde{C} > 0$. 

Note that in the absence of an explicit constant $\widetilde{C}$, this inequality is only interesting in the regime that $\Var(S) \to \infty$, where central limit theorems are often more viable.  In particular, it yields trivial results in the Poisson regime, namely $M(S) \leq O(1)$, when the variance of $S$ can be bounded away from zero. In regimes where the variance tends to $0$, the result is of an order worse than the trivial $M(S) \leq 1$. Whereas, our bounds derived in Section \ref{2} provide meaningful quantitative estimates in all regimes.



In more recent work, explicit constants have been obtained for a variant of the Miroshnikov-Rogozin inequality. To this end, for $\theta$ a unit vector in $\mathbb{R}^n$ define
\[
    p_\theta(t) = \vol_{n-1} \left\{ x \in \R^n : \|x\|_\infty \leq \frac 1 2, \langle x, \theta \rangle = t \right\}.
\]



\begin{thm}[Bobkov-Chistyakov \cite{BC}]
    Given $\lambda \geq \left(\sum_{k=1}^n \lambda_k^2 \right)^{\frac 1 2}$ and $X_k$ independent random variables with $S = X_1 + \cdots + X_n$,
    \[
        Q(S; \lambda) \leq  \frac{ 2^{\frac 3 2} \lambda}{c}   \left( \sum_{k=1}^n \lambda_k^2 Q^{-2}(X_k; \lambda_k) \right)^{- \frac 1 2},
    \]
    where  $$c \coloneqq  \inf_{|t| <\frac 1 2, \|\theta\| = 1 } p_\theta(t).$$ 
\end{thm}
Moreover, Bobkov and Chistyakov showed that $c \geq 0.00095$ independent of dimension.
Later Melbourne, Tkocz, and Wyczesany \cite{MTW} showed that the body $\{ x \in \R^n : \|x\|_\infty \leq \frac 1 2 \}$ can be replaced by any isotropic convex body $K$.  More precisely for a convex body $K$ such that  $\int_K x_i x_j dx = L_K^2 \delta_{ij}$,  for some constant $L_K > 0$, 
\[
    \inf_{\theta, |t| < L_K \sqrt{3}}  \vol_{n-1} \left \{ x \in K : \langle x , \theta \rangle = t \right \} \geq \frac{e^{- \sqrt{6}}}{\sqrt{2 L_K^2}}.
\]
This result is proven sharp for high dimensional double cones, but even in the case of the cube where $L_K = \frac{1}{\sqrt{12}}$ it gives $c \geq \sqrt{6} e^{- \sqrt{6}}$, so that Bobkov and Chistyakov can be stated as 

    \[
        Q(S; \lambda) \leq  \frac{2 e^{\sqrt{6}}}{\sqrt{3}} \lambda  \left( \sum_{k=1}^n \lambda_k^2 Q^{-2}(X_k; \lambda_k) \right)^{- \frac 1 2}.
    \]
For reference, numerically $\frac{2 e^{\sqrt{6}}}{\sqrt{3}} \ \approx 13.3742$. Hence for $C =\frac{2 e^{\sqrt{6}}}{\sqrt{3}} $, $X_k$ integer valued and $\sum_{k=1}^n \lambda_k^2 \leq \lambda^2 < 1$, by Bobkov-Chistyakov
\[
    M(S) = Q(S; \lambda) \leq C \lambda \left( \sum_{k=1}^n \lambda_k^2 Q^{-2}(X_k; \lambda_k ) \right)^{- \frac 1 2} = C \lambda \left( \sum_{k=1}^n \lambda_k^2 M^{-2}(X_k )\right)^{- \frac 1 2}.
\]
Minimizing the right-hand side over $\lambda_k$, that is when $\lambda_j = \lambda$ for $j$ such that $M^{-2}(X_j) = \max_{1 \leq k \leq n} M^{-2}(X_k)$ and $\lambda_l = 0$ for $l \neq j$, gives
\[
    M(S) \leq C \left( \max_{1 \leq k \leq n} M^{-2}(X_k) \right)^{- \frac 1 2} = C \min_{1 \leq k \leq n} M(X_k).
\]
However, this is a trivial bound as Young's convolution inequality yields 
$$ M(S) \leq \min_{1 \leq k \leq n} M(X_k). $$







\section{A Littlewood-Offord type problem for arithmetic progressions}\label{4}

Given independent $\mathbb{Z}$-valued random variables $X_1, \dots, X_n$, $n \geq 1$, and $a = (a_1, \dots, a_n) \in (\mathbb{R} \setminus \{0\} )^n$, we ask for an upper bound on 
\[
    \sup_{x \in \R} \, \mathbb{P}(a \cdot X \in A_{l,m}(x))
\]
where $A_{l,m}(x)$ is an arithmetic progression of length $l \geq 1$, that is $A_{l,m}(x) = \{x + mj \}_{j=1}^l$ for $m, x \in \mathbb{R}$.  In the case that $l = 1$ this corresponds to the classical Littlewood-Offord problem for the variables $X_k$.

\begin{proof}[Proof of Theorem \ref{main3}]
Let $Y$ be a discrete random variable independent of $U_l$, where $U_l$ is uniform on $\{1,2,\dots, l\}$. Then,
\begin{align*}
    \mathbb{P}(Y - m U_l = x) 
        &= 
            \sum_{k=1}^l \mathbb{P}(U_l = k, Y = x + mk ) 
                \\
        &=
            \frac 1 l \sum_{k=1}^l \mathbb{P}(Y = x +mk)
                \\
        &=
            \frac 1 l \mathbb{P}( Y \in A_{l,m}(x)).
\end{align*}
Thus, this Littlewood-Offord problem for arithmetic progressions is equivalent to finding upper bounds on $M(a \cdot X - m U_l)$. When the $X_k$'s are discrete log-concave, one may thus apply Theorem \ref{main1} to obtain
$$ \sup_{x \in \R} \, \mathbb{P}(a \cdot X \in A_{l,m}(x)) = l \, M(a \cdot X - m U_l) \leq \frac{l}{\sqrt{1 + c \left( \sum_{k=1}^n \Var(X_i) + \Var(U_l) \right)}}, $$
which is the desired result since $\Var(U_l) = (l^2-1)/12$.
\end{proof}



\begin{remark}
Let us note that Theorem \ref{main3} is sharp up to an absolute constant when taking supremum over all $m \neq 0$, as
\begin{equation*}
\sup_{m \neq 0} \sup_{a \in (\R \setminus \{0\})^n} \sup_{x \in \R} \, \Prob(a_1 X_1 + \cdots + a_n X_n \in A_{l,m}(x)) \geq \frac{l}{\sqrt{1 + 12 \sum_{k=1}^n \Var(X_k) + l^2-1}},
\end{equation*}
due to the lower bound in \eqref{M-bound}.
\end{remark}


\section{Specific case of the Bernoulli and uniform distributions}

\subsection{Bernoulli distribution}



This section focuses on strengthening Theorem \ref{main3} for the Bernoulli distribution. Let $U_l$ be uniform on $\{1, \dots, l\}$. It has been shown in \cite[Proof of Theorem 1.7]{MR} that for all $p \geq 2$ the Fourier transform of $U_l$ satisfies
\[
    \int_{- \frac 1 2}^{\frac 1 2} \left| \mathbb{E} e^{2 \pi i t U_l} \right|^p dt \leq \int_{- A}^{A} e^{- \pi (l^2 - 1) p t^2 /2} dt,
\]
where $A$ is determined implicitly through the equation
\begin{equation}\label{A-def}
    \int_{-A}^A e^{- \pi (l^2 - 1) t^2} dt = \int_{-\frac 1 2 }^{\frac 12} \left| \mathbb{E}e^{2 \pi i t U_l} \right|^2 dt = \frac 1 l.
\end{equation}
The first equality in \eqref{A-def} gives the implicit definition of $A$ (in terms of $l$), while the second equality comes from recalling that $U_l$ is uniform and using Parseval identity. Therefore,
\begin{equation}\label{A-1} \frac 1 {2\pi} \int_{-\pi}^{\pi} \left| \mathbb{E} e^{itU_l} \right|^p dt \leq 2  \int_{0}^{A} e^{- \pi (l^2 - 1)p t^2 /2} dt = \frac{2}{\sqrt{ \pi  (l^2-1) p}} \int_0^{\sqrt{\pi (l^2-1)A^2 p}} e^{-x^2/2} dx = 2A \Phi (c p),
\end{equation}
where $\Phi(x) \coloneqq \frac 1 {\sqrt{x}} \int_0^{\sqrt{x}} e^{-t^2/2} dt$ and $c = \pi (l^2 -1) A^2$. Note that 
\[
A = \frac{1}{\sqrt{ \pi(l^2 - 1) }} \cdot \operatorname{erf}^{-1}\left( \frac{\sqrt{l^2 - 1}}{l} \right),
\]
where $\operatorname{erf}(x) = \frac{2}{\sqrt{\pi}} \int_0^{x} e^{-t^2} dt$ is the error function. On the other hand, it has been shown in \cite[Theorem 2.8]{MMR} that if $X$ is a Bernoulli random variable with variance $\sigma^2$ and $q \geq 1$,
\begin{equation}\label{A-2}
        \frac 1 {2\pi} \int_{-\pi}^{\pi} \left| \mathbb{E} e^{itX} \right|^q dt \leq \frac 1 {\sqrt{6 \sigma^2 q}} \int_0^{\sqrt{6 \sigma^2 q}} e^{-t^2/2} = \Phi(6 \sigma^2 q).
\end{equation}
    
Let us now consider $X_1, \dots, X_n$ independent Bernoulli random variables with variance $\sigma_k^2$ and $U_l$ uniform on $\{1, \dots, l\}$, and denote $v=(\pm1, \dots, \pm 1)$ and $v_0=\pm 1$ any choice of signs, then
\begin{eqnarray*}
\|f_{\sum_{k=1}^n v_k X_k + v_0 U_l}\|_{\infty} & \leq & \|\widehat{f}_{\sum_{k=1}^n v_k X_k + v_0 U_l}\|_{1} \\ & = & \frac 1 { 2\pi } \int_{- \pi}^{\pi} \left| \mathbb{E} e^{i t (v \cdot X  + v_0 U_l)} \right| dt \\ & \leq & \frac{1}{2 \pi} \left(\int_{- \pi}^{\pi} \left|\mathbb{E} e^{it U_l}\right|^p dt \right)^{\frac 1 p} \prod_{k=1}^n \left(\int_{-\pi}^{\pi} \left| \mathbb{E} e^{it X_k} \right|^{q_k} dt \right)^{\frac 1 {q_k}} \\ & \leq &  \left( 2A \Phi (\pi (l^2 -1) A^2 p) \right)^{\frac{1}{p}} \, \prod_{k=1}^n \Phi \left( 6 \sigma_k^2 q_k \right)^{\frac{1}{q_k}},
\end{eqnarray*}
where we have used the Hausdorff-Young inequality, H\"older's inequality with $\frac 1 p + \sum_{k=1}^n \frac 1 {q_k} = 1$,  the independence of the variables, and the bounds \eqref{A-1}, \eqref{A-2}. Choosing $C = \pi (l^2-1)A^2 + 6 \sum_{k=1}^n \sigma_k^2$, and setting $q_k = \frac{C}{6 \sigma_k^2}  $ and $p = \frac{C}{\pi (l^2-1) A^2}$, then if $p \geq 2$ we have
\begin{align*}
        \sup_{x \in \R} \mathbb{P}(a \cdot X \in A_l)
            &\leq 
                l \, \|f_{\sum v_i X_i + v_0 U_l}\|_{\infty}
                    \\
            & \leq l (2A)^{\frac{1}{p}} \Phi (C) ^{\frac{1}{p}} \, \prod_{k=1}^n \Phi \left( C \right)^{\frac{1}{q_k}}\\
            &=
                l (2A)^{\frac{1}{p}} \Phi(C).
\end{align*}
Using the bound $\Phi(z) \leq \frac{1}{\sqrt{1+\frac{z}{3}}}$ holding for all $z > 0$ (see \cite[Lemma 2.9]{MMR})
we deduce
\begin{equation}\label{Bern}
\sup_{x \in \R} \mathbb{P}(a \cdot X \in A_l) \leq  (2A)^{\frac{1}{p}} \frac{l}{\sqrt{1 + 2 \sum_{k=1}^n \Var(X_k) + \frac{l^2-1}{12} 4 \pi A^2} }.
\end{equation}




\begin{remark}
    
Let us compare the bound \eqref{Bern} with theorem \ref{main3}. Since $2A \leq 1$, we have 
$$ (2A)^{\frac{1}{p}} \frac{l}{\sqrt{1 + 2 \sum_{k=1}^n \Var(X_k) + \frac{l^2-1}{12} 4 \pi A^2} } \leq \frac{l}{\sqrt{1 + 2 \sum_{k=1}^n \Var(X_k) + \frac{l^2-1}{12} 4 \pi A^2} }. $$

Therefore, for $l=1$, we recover the bound \eqref{answer-2} proved in \cite{MMR}. Moreover, for $l = 2$, one can check numerically that $4 \pi A^2 \geq 1$, so that the bound \eqref{Bern} is always better for Bernoulli distribution than Theorem \ref{main3}. Note that for fixed length $l$, the bound \eqref{Bern} is stronger as the variance grows.

However, the bound is not always applicable when $l \geq 2$, as $p = 1 + \frac{6 \sum_{k=1}^n \Var(X_k)}{\pi (l^2-1) A^2}$ needs to be greater than or equal to 2. Hence, one may not choose variances that are too small compared to the length $l$.

\end{remark}





\subsection{Discrete uniform distribution}

This section specializes to the uniform distribution. In particular, we recover and extend a result of Madiman and Woo \cite{MW} on the entropy power inequality for discrete uniform distributions.

Let $U$ be a uniform distribution on consecutive integers $\{a, \dots, b\}$, with $a,b \in \Z$. Denote $l = b-a+1$. Note that $\Var(U) = (l^2-1)/12$, therefore its Fourier transform satisfies
$$ \|\widehat{f_U}\|_2^2 = \frac{1}{l} = \frac{1}{\sqrt{1 + 12 \Var(U)}}, $$
where the first identity follows from Parseval. Therefore, using \cite[Lemma 2.7]{MMR}, we obtain for $n$ independent uniformly distributed random variables $U_1, \dots, U_n$
$$ \|\widehat{f}_{\sum_{k=1}^n U_k}\|_2^2 \leq \frac{1}{\sqrt{1 + 12 \sum_{k=1}^n \Var(U_k)}}. $$
Using the Parseval identity, this leads to
\begin{equation*}
H_2 \left( \sum_{k=1}^n U_k \right) = \log(\|p_{\sum_{k=1}^n U_k}\|_2^{-2}) = \log(\|\widehat{f}_{\sum_{k=1}^n U_k} \|_2^{-2}) \geq \frac{1}{2} \log \left( 1 + 12 \sum_{k=1}^n \Var(U_k) \right).
\end{equation*}
By monotonicity of entropy, we deduce that for all $\alpha \leq 2$,
\begin{equation}\label{unif}
H_{\alpha} \left( \sum_{k=1}^n U_k \right) \geq \frac{1}{2} \log \left( 1 + 12 \sum_{k=1}^n \Var(U_k) \right).
\end{equation}
We are now ready to prove Theorem \ref{EPI-unif}.

\begin{proof}[Proof of Theorem \ref{EPI-unif}]
    Let $\alpha \leq 2$. Denoting $\Delta_\alpha(X) = N_\alpha(X) - 1$, which reflects the variance better than the entropy power for discrete distributions, we deduce from \eqref{unif} that
\[
    \Delta_\alpha \left( \sum_{k=1}^n U_k \right) \geq 12 \sum_{k=1}^n \Var(U_k).
\]
However, we note that for a uniform random variable $U$ on an integer interval, $\Delta_\alpha(U) = 12 \ \Var(U)$, for any $\alpha$. Thus we have
\begin{equation}\label{bound-unif}
    \Delta_\alpha \left( \sum_{k=1}^n U_k \right) \geq \sum_{k=1}^n \Delta_\alpha(U_k),
\end{equation}
for $\alpha \leq 2$. Moreover, for $X_k$ uniform on any set of integers, $X_k^{\#}$ has a uniform distribution on an integer interval and hence is log-concave, thus $X_k^{\#} \sim U_k$ for $U_k$ uniform on an integer interval, therefore
\begin{equation*}
\Delta_\alpha \left( \sum_{k=1}^n X_k \right) \geq \Delta_\alpha \left( \sum_{k=1}^n X_k^{\#} \right) \geq \sum_{k=1}^n \Delta_\alpha \left( X_k^{\#} \right) = \sum_{k=1}^n \Delta_\alpha(X_k),
\end{equation*}
where the first inequality comes from Theorem \ref{rearrance-2} together with Schur concavity of R\'enyi entropy, and the second inequality from \eqref{bound-unif}.
\end{proof}


\begin{remark}
    \begin{enumerate}
        \item Let us note that \eqref{bound-unif} implies Theorem \ref{main2} with $c=12$ and $\alpha \in [0,2]$ when the random variables are uniformly distributed.

        \item Employing the relation
$$ H_{\alpha}(X) \leq H_{\infty}(X) + \log(\alpha^{\frac{1}{\alpha - 1}}) $$
obtained in \cite{MT}, which is valid for all log-concave distributions and $0 < \alpha < \infty$, we deduce by taking $\alpha =2$ that
$$ N_{\infty} \left( \sum_{k=1}^n U_k \right) \geq \frac{1}{4} N_{2} \left( \sum_{k=1}^n U_k \right) \geq \frac{1}{4} + 3  \sum_{k=1}^n \Var(U_k), $$
where the second inequality comes from \eqref{unif}. Equivalently,
$$ M \left( \sum_{k=1}^n U_k \right) \leq \frac{1}{\sqrt{\frac{1}{4} + 3  \sum_{k=1}^n \Var(U_k)}}, $$
which is an improvement of Theorem \ref{main1} whenever the random variables are uniformly distributed on at least 3 points.
    \end{enumerate}
\end{remark}









\begin{thebibliography}{10}


\bibitem{A} H. Aravinda, Entropy-variance inequalities for discrete log-concave random variables via degree of
freedom. Discrete Mathematics, 347(1):113683, 2024.

\bibitem{BFK} A. Bandeira, A. Ferber, M. Kwan. Resilience for the {L}ittlewood–{O}fford problem. Advances in Mathematics, 319:292–312, 10 2017.

\bibitem{BC} Bobkov, S., Chistyakov, G., On concentration functions of random variables. J. Theoret. Probab. 28 (2015), no. 3, 976--988.

\bibitem{BM} S. Bobkov, A. Marsiglietti, Entropic CLT for smoothed convolutions and associated entropy bounds, Int. Math. Res. Not. IMRN 2020, no. 21, 8057--8080.

\bibitem{BMM} S. Bobkov, A. Marsiglietti, J. Melbourne. Concentration functions and entropy bounds for discrete log-concave distributions, Combin. Probab. Comput. 31 (2022), no. 1, 54--72.

\bibitem{Bra}  P. Br\"and\'en. Unimodality, log-concavity, real-rootedness and beyond. Handbook of enumerative combinatorics, Discrete Math. Appl., CRC Press:437--483, 2015.

\bibitem{Bre} F. Brenti. Log-concave and unimodal sequences in algebra, combinatorics, and geometry: an update. In Jerusalem combinatorics ’93, volume 178 of Contemp. Math., pages 71--89.
Amer. Math. Soc., Providence, RI, 1994.

\bibitem{E} P. Erd\"os, On a lemma of Littlewood and Offord, Bulletin of the American Mathematical Society 51 (1945), no. 12, 898–902.

\bibitem{FKS} J. Fox, M. Kwan, L. Sauermann, Combinatorial anti-concentration inequalities, with applications, Math. Proc. Cambridge Philos. Soc. 171 (2021), no. 2, 227--248.

\bibitem{GMPP} T. Gaxiola, J. Melbourne, V. Pigno, E. Pollard, Anti-concentration inequalities for log-concave variables on the real line. Preprint, arXiv:2505.05793.

\bibitem{H} G. Hal\'asz, Estimates for the concentration function of combinatorial number theory and probability, Periodica Mathematica Hungarica 8 (1977), 197--211.

\bibitem{JMNS} J. Jakimiuk, D. Murawski, P. Nayar, S. S\l obodianiuk. Log-concavity and discrete degrees of freedom. Discrete Mathematics, 347(6):114020, 2024.

\bibitem{JG} O. Johnson, C. Goldschmidt, Preservation of log-concavity on summation. ESAIM Probab. Stat. 10 (2006), 206--215 (electronic).

\bibitem{LO} J. E. Littlewood, A. C. Offord, On the number of real roots of a random algebraic equation. III, Math. Mat. Sbornik N.S. 12 (1943), no. 3, 277–286.

\bibitem{JK} Juškevičius, T., Kurauskas, V. (2021). On the Littlewood‐Offord problem for arbitrary distributions. Random Structures \& Algorithms, 58(2), 370-380.

\bibitem{K1} D. Kleitman, On a lemma of Littlewood and Offord on the distribution of certain sums, Math Z 90, 251–259 (1965).

\bibitem{K2} D. Kleitman, On a Lemma of Littlewood and Offord on the Distributions of Linear Combinations of Vectors, Advances in Mathematics, 5, 155--157 (1970).

\bibitem{MMR} M. Madiman, J. Melbourne, C. Roberto, Bernoulli sums and Rényi entropy inequalities, Bernoulli 29 (2023), no. 2, 1578–1599.

\bibitem{MWW} M. Madiman, L. Wang, J. O. Woo, Majorization and R\'enyi Entropy Inequalities via Sperner Theory, Discrete Mathematics (AEGT 2017 Special issue edited by S. Cioaba, R. Coulter, E. Fiorini, Q. Xiang, F. Pfender), vol. 342, no. 10, pp. 2911--2923.

\bibitem{MOA} A. W. Marshall, I. Olkin, B. C. Arnold, Inequalities: theory of majorization and its applications, 2nd Edition, Springer Series in Statistics, Springer, New York, 2011.

\bibitem{Massey} J. L. Massey, On the entropy of integer-valued random variables. In: Proc. 1988 Beijing Int. Workshop on
Information Theory, pages C1.1-C1.4, July 1988.

\bibitem{MR} J. Melbourne, C. Roberto, Transport-majorization to analytic and geometric inequalities. JFA

\bibitem{MT} J. Melbourne, T. Tkocz, Reversal of R\'enyi entropy inequalities under log-concavity. IEEE Trans. Inform. Theory, 67 (2020), no. 1, 45--51.

\bibitem{MTW} J. Melbourne, T. Tkocz, K. Wyczesany, A Rényi entropy interpretation of anti-concentration and noncentral sections of convex bodies. Preprint, arXiv:2406.04200 [math.PR].


\bibitem{M-R} Mirosnikov, A. L., Rogozin, B. A., Inequalities for concentration functions. Teor. Veroyatnost. i Primenen. 25 (1980), no. 1, 178--183.

\bibitem{NV} H. Nguyen, V. Vu, Optimal inverse Littlewood-Offord theorems, Advances in Mathematics 226 (2011), no. 6, 5298--5319.

\bibitem{SW} A. Saumard, J. A. Wellner. Log-concavity and strong log-concavity: a review. Statistics surveys, 8:45, 2014.

\bibitem{S} M. Singhal, Erd\"os-Littlewood-Offord problem with arbitrary probabilities, Discrete Math. 345 (2022), no. 11, Paper No. 113005, 13 pp.

\bibitem{Sta} R. P. Stanley, Log-concave and unimodal sequences in algebra, combinatorics, and geometry. In Graph theory and its applications: East and West (Jinan, 1986), volume 576 of Ann. New York Acad. Sci., pages 500--535. New York Acad. Sci., New York, 1989.

\bibitem{TV} T. Tao, V. H. Vu, A sharp inverse Littlewood-Offord theorem, Random Structures \& Algorithms 37 (2010), no. 4, 525-539.

\bibitem{MW} J. O. Woo, M. Madiman, A discrete entropy power inequality for uniform distributions, Proc. IEEE Intl. Symp. Inform. Theory, Hong Kong, China, 2015.

\end{thebibliography}




\iffalse


\vskip1cm

\noindent Arnaud Marsiglietti \\
Department of Mathematics \\
University of Florida \\
Gainesville, FL 32611, USA \\
E-mail: a.marsiglietti@ufl.edu

\vspace{0.8cm}

\noindent James Melbourne \\
Probabilidad y Estad\'isticas\\
Centro de Investigaci\'ones en Matem\'aticas \\
Guanajuato, GTO 36023, MX \\
E-mail: james.melbourne@cimat.mx


\fi




\end{document}
